\documentclass[11pt]{article}
\usepackage[a4paper,margin=26mm]{geometry}
\usepackage{amsmath,amssymb,amsthm,booktabs,array}
\usepackage[T1]{fontenc}
\usepackage{lmodern,microtype}
\usepackage[hidelinks]{hyperref}
\hypersetup{pdftitle={Finite causal spin-two transfer and partial refinement in Lorentzian Regge calculus},pdfauthor={}}
\newtheorem{theorem}{Theorem}[section]
\newtheorem{proposition}[theorem]{Proposition}
\newtheorem{corollary}[theorem]{Corollary}
\theoremstyle{remark}\newtheorem{remark}[theorem]{Remark}
\DeclareMathOperator{\rank}{rank}
\DeclareMathOperator{\diag}{diag}
\DeclareMathOperator{\tr}{tr}
\newcommand{\R}{\mathbb R}
\newcommand{\Id}{\mathbb I}
\title{Finite causal spin-two transfer and partial refinement in Lorentzian Regge calculus}
\author{}
\date{29 October 2026}
\begin{document}
\maketitle
\begin{abstract}
A finite Lorentzian Regge region can support an angular spin-two transfer channel without defining a continuum graviton.
We give an explicit cone over a longitudinal polygon joined to a transverse hexagon, specify the complete complex action and its physical radial Schur reduction, and isolate the real two-dimensional angular-weight-two representation of the dihedral group.
The symmetric coarse witness has two approximately degenerate elliptic transfer phases near $0.815035$.
An exact representation-and-support argument bounds its cap-to-cap mixed Hessian rank by two at every nondegenerate level of this architecture: only one internal spin-two pair mediates transmission.
Partial longitudinal subdivisions retain the pair, but the fifth tested level has reciprocal-real multipliers near $1.29025$ and $0.77504$ in the stated reduced endpoint convention.
We distinguish this finite dependence from an asymptotic obstruction: fixed transverse edges prevent the maximum mesh size from vanishing, and consecutive integer subdivisions are not all nested.
We also correct a spectral-invariance claim: independent endpoint canonical frames act by left--right equivalence, whereas a tracked endpoint identification defines a conjugacy-invariant return map.
An explicit scalar counterterm and common dilation match the selected quadratic channel to a chosen target kernel.
This is a completed finite calibration, not a local perfect Regge action.
The accompanying portable calculation proves the decisive algebra and verifies frozen complete-action outputs; positive energy, retarded propagation and a full continuum limit remain open.
\end{abstract}

\section{Question and relation to established results}
\label{sec:intro}
Canonical evolution in Regge calculus can change the number of variables on a spatial boundary.
Hamilton's principal function gives a natural canonical relation, including the pre- and post-constraints associated with this change~\cite{DittrichHoehn2012,Hoehn2014}.
In linearized four-dimensional gravity, a $2\to3$ spatial Pachner move creates a lattice-graviton variable, while a $3\to2$ move removes one and introduces a nontrivial equation~\cite{Hoehn2015}.
Neither a finite mode count nor symplecticity alone identifies the two helicities of a continuum massless field.

The contrast between four-dimensional $4$--$2$ and $3$--$3$ bulk moves is likewise established.
On the flat branch the former eliminates an internal edge, whereas the latter contains a boundary-determined curvature hinge and generally changes the Hamilton--Jacobi function~\cite{BorissovaDittrich2023}.
Our subject is a larger explicit Lorentzian witness in which causal boundary geometry, an angular representation, complete-action dynamics and partial refinement can be assessed together.
We do not claim new canonical Pachner theory, a first lattice graviton, or a general failure of Regge calculus.

The concrete questions are whether the finite region transmits a dihedral angular-weight-two pair, what fixes the observed rank, and what a comparison across its subdivisions actually means.
The answer contains positive finite content and specific limitations.
The coarse transfer is elliptic to differentiation accuracy, the architecture enforces a rank bound, and the fifth selected transfer is hyperbolic in its stated convention.
The geometry is only partially refined, however, and a spectrum of a map between distinct endpoint spaces needs an identification before it is coordinate invariant.
Improved actions and dynamical embedding maps are established tools for addressing the wider problem~\cite{BahrDittrich2009,Dittrich2012}; the reduced calibration below solves only the scalar channel-matching question.

\section{An explicit Lorentzian region}
\label{sec:geometry}
All coordinates and the overall action normalization are dimensionless in the numerical witness.
Use $\eta=\diag(-1,1,1,1)$ in coordinates $(t,x,y,z)$, an internal vertex $o=0$, and five longitudinal points $q_0,\ldots,q_4$ in the $t$--$z$ plane.
Their nonzero coordinates are
\begin{equation}
\begin{array}{c|rr}
 &t&z\\\hline
q_0&0&-1.5709397991101974\\
q_1&0.3762757793700152&-0.42646145218382514\\
q_2&0.7702937141192545& 1.1723475896580613\\
q_3&-0.7702937141192545&1.1723475896580613\\
q_4&-0.3762757793700152&-0.42646145218382514
\end{array}.
\label{eq:qpoints}
\end{equation}
The transverse points are
\begin{equation}
r_j=(0,\cos\phi_j,\sin\phi_j,0),\qquad
\phi_j=j\pi/3,\quad j=0,\ldots,5.
\label{eq:rpoints}
\end{equation}
The cyclic orders are $q_0q_1q_2q_3q_4q_0$ and $r_0\cdots r_5r_0$.
Cone their join from $o$: for every longitudinal cycle edge $[q,q']$ and transverse edge $[r_j,r_{j+1}]$, include
\begin{equation}
\sigma=[o,q,q',r_j,r_{j+1}].
\label{eq:simplices}
\end{equation}
The join of the two polygonal circles is a three-sphere; its cone is a combinatorial four-ball.
Its boundary tetrahedra are exactly $[q,q',r_j,r_{j+1}]$.

At subdivision level $m$, insert $(1-j/m)q_0+(j/m)q_1$ on $[q_0,q_1]$ and $(1-j/m)q_4+(j/m)q_0$ on $[q_4,q_0]$, for $j=1,\ldots,m-1$.
The longitudinal cycle has $2m+3$ vertices, and the bulk has $6(2m+3)$ four-simplices.
Time reflection exchanges the two subdivided edges and $q_1\leftrightarrow q_4$, $q_2\leftrightarrow q_3$.

For a simplex with base vertex $0$ and signed squared lengths $s_{ab}$, the Gram matrix is
\begin{equation}
G_{ab}=\frac{s_{0a}+s_{0b}-s_{ab}}2.
\label{eq:gram}
\end{equation}
Every tested simplex has one negative and three positive eigenvalues.
For a boundary tetrahedron with outward normal $n$, $n^T\eta n<0$ means it is spacelike, and the sign of $n^0$ distinguishes past from future.
The remaining boundary strip is timelike.
Table~\ref{tab:geometry} gives direct checks on the specified coordinates; no normal is null.

\begin{table}[htbp]
\centering\small
\caption{Geometry of the partial-subdivision family. The last two columns are minimum Gram eigenvalue margins across the four-simplices, using Eq.~\eqref{eq:gram} based at $o$.}
\label{tab:geometry}
\begin{tabular}{rrrrr}
\toprule
$m$&simplices&past/future/side&$\min|\lambda_-|$&$\min\lambda_+$\\
\midrule
1&30&12/12/6&0.132328&0.500000\\
2&42&18/18/6&0.0252817&0.500000\\
3&54&24/24/6&0.0100139&0.500000\\
4&66&30/30/6&0.00530682&0.500000\\
5&78&36/36/6&0.00327578&0.413006\\
\bottomrule
\end{tabular}
\end{table}

Fix the endpoint cones and side strip.
The initial cap is the subdivided path from $q_0$ through $q_1$ to $q_2$, joined to the transverse circle; the final cap is its reflected path from $q_3$ through $q_4$ to $q_0$.
Its free longitudinal vertices are the $m$ non-endpoint sites on each path.
Relative cap edges comprise its $m+1$ path edges and the $6m$ edges from those sites to the hexagon.
Thus there are $7m+1$ relative squared lengths per cap.
The geometry and normals specify a time-oriented boundary region; they do not by themselves prove retarded response.

\section{Complete action, branch and physical quotient}
\label{sec:action}
We use the complexified Lorentzian Regge action including all boundary hinges~\cite{BorissovaDittrich2023}.
For a triangle $h=(a,b,c)$, its signed area square and area are
\begin{equation}
V_h=\frac{4s_{ab}s_{ac}-(s_{ab}+s_{ac}-s_{bc})^2}{16},
\qquad A_h=\sqrt{V_h},
\label{eq:area}
\end{equation}
with the principal square root, positive imaginary on the negative real axis.
The producer realizes each off-flat simplex locally from Eq.~\eqref{eq:gram} by its spectral decomposition, so no global off-flat coordinate embedding is assumed.

For the two tetrahedral facets meeting at $h$, let $n_1$ be outward and $n_2$ be the negative of the other outward normal.
Set $d=n_1\cdot n_2$ and $D=(n_1\cdot n_1)(n_2\cdot n_2)-d^2$.
The branch used by this mixed-facet witness is
\begin{equation}
\theta_{\sigma h}=-i\log^-
\frac{d+\rho_h i\sqrt{D}}{
\sqrt{n_1\cdot n_1}\sqrt{n_2\cdot n_2}},\qquad
\rho_h=\begin{cases}+1,&A_h\text{ imaginary},\\-1,&A_h\text{ real}.
\end{cases}
\label{eq:angle}
\end{equation}
Here $\log^-$ takes argument $-\pi$ on the negative real cut.
This is the local mixed-causal-sector evaluator of the retained computation, not a prescription for crossing arbitrary branch singularities.
The recorded producer checks supporting this sheet are inherited; this paper's inexpensive replay does not perform a new interval proof of branch regularity.

Writing $\Theta_h=\sum_{\sigma\supset h}\theta_{\sigma h}$, the implementation selects a locally constant integer
\begin{equation}
p_h=\begin{cases}2\pi,&h\text{ internal},\\\pi,&h\text{ boundary},\end{cases}
\quad k_h=\operatorname{nint}(-\Re\Theta_h/p_h),\quad
\epsilon_h=k_hp_h+\Theta_h,
\quad S=-i\sum_h A_h\epsilon_h.
\label{eq:action}
\end{equation}
The branch integer is held effectively constant inside a regular local component; a universal boundary integer must not be substituted.
The sum includes internal and boundary triangles and uses the declared flat-background sheet.

An embedded displacement acts on squared lengths as
\begin{equation}
\delta s_{ab}=2(x_a-x_b)^T\eta(\delta x_a-\delta x_b).
\label{eq:tangent}
\end{equation}
On either cap this tangent has rank $4m$ for the declared moving sites, leaving $3m+1$ quotient directions.
An orthonormal basis of the left-null complement represents that quotient.
For the internal radial edges $[o,v]$, remove the four directions generated by moving $o$; their physical complement has dimension $2m+5$.
These are linearized flat-background quotients, not a claim that every curved Regge configuration has exact vertex-displacement gauge symmetry.

Let $z$ denote these physical internal coordinates and let $x,y$ be cap coordinates.
Differentiating Eq.~\eqref{eq:action} produces a symmetric, generally complex Hessian $H$.
If $H_{zz}$ is invertible, the quadratic Hamilton principal function has boundary Hessian
\begin{equation}
K=H_{\partial\partial}-H_{\partial z}H_{zz}^{-1}H_{z\partial}.
\label{eq:schur}
\end{equation}
The inverse is taken only after removing the four central displacement directions.
Equivalently, the physical internal solution is $z=-H_{zz}^{-1}H_{z\partial}(x,y)$.
The true linearized internal Regge equations are area-moment equations,
\begin{equation}
\delta E_e=\sum_{h\supset e}
\left.\frac{\partial A_h}{\partial s_e}\right|_0
\delta\epsilon_h=0.
\label{eq:reggeeq}
\end{equation}
An arbitrary tensor-screen nullity is not a substitute for Eq.~\eqref{eq:reggeeq} or for a helicity test.

The retained calculations use central finite differences followed by Richardson extrapolation,
\begin{equation}
H^{\mathrm R}(h)=\frac{4H^{\mathrm{cd}}(h)-H^{\mathrm{cd}}(2h)}3,
\qquad h=1.2\times10^{-4}
\label{eq:richardson}
\end{equation}
in the normalized quotient coordinates for the two freshly reproduced levels.
Residual imaginary parts and nondegenerate-pair splittings are retained as numerical diagnostics rather than discarded before differentiation.

\section{The angular sector and its rank bound}
\label{sec:rank}
At a longitudinal site use raw relative-edge variations proportional to
\begin{equation}
u_2(\phi_j)=(\cos 2\phi_j,\sin 2\phi_j)
\label{eq:e2}
\end{equation}
on the edges to $r_j$, and zero on the longitudinal path edges.
Projection into the cap quotient followed by orthonormalization gives $2m$ spin-two samples per cap.
The coarse raw samples lie in the quotient to a relative residual below $4.4\times10^{-16}$ in the retained replay.

On a single real pair, the transverse rotation and reflection are
\begin{equation}
R=\begin{pmatrix}-1/2&-\sqrt3/2\\\sqrt3/2&-1/2\end{pmatrix},
\qquad F=\diag(1,-1).
\label{eq:representation}
\end{equation}
Solving $R^TXR=X$ gives $X=\alpha\Id+\beta J_2$, where
$J_2=\left(\begin{smallmatrix}0&-1\\1&0\end{smallmatrix}\right)$.
Imposing $F^TXF=X$ gives $X=\alpha\Id$.
Consequently an exactly invariant Hessian is diagonal between inequivalent angular sectors, and every block between two copies of this real dihedral representation is a scalar times $\Id_2$.
This is angular weight two modulo six: six angular samples cannot uniquely distinguish it from all continuum weights that alias to the same finite representation.

\begin{theorem}[Single internal angular channel]
\label{thm:rank}
For the cone in Eqs.~\eqref{eq:qpoints}--\eqref{eq:simplices}, at any subdivision level on a dihedral-invariant, differentiable branch, assume the physical internal spin-two Hessian is nonsingular.
The cap-to-cap mixed spin-two Schur block has the form
\begin{equation}
B_m=-\frac{uv^T}{k}\otimes\Id_2,
\label{eq:bottleneck}
\end{equation}
for longitudinal coupling vectors $u,v\in\R^m$ and a nonzero scalar $k$ in the real quadratic idealization.
Its rank is at most two, and is two if both coupling vectors are nonzero.
The same algebraic rank statement holds over $\mathbb C$ for the analytic complex Hessian.
\end{theorem}
\begin{proof}
Every radial edge from $o$ to a longitudinal vertex is invariant under the transverse dihedral group.
The six radial edges to the transverse hexagon form its vertex permutation representation.
Their Fourier decomposition contains precisely one real angular-weight-two pair.
The four central displacement directions transform as two scalars and one transverse vector pair, so removing them does not remove or add a spin-two copy.
Thus the complete physical internal spin-two space is one pair.
Dihedral invariance makes its self block $k\Id_2$ and the couplings to cap multiplicity spaces $u\otimes\Id_2$ and $v\otimes\Id_2$.

No simplex contains a moving longitudinal vertex of each cap: the two cap paths meet only at their fixed endpoint, and the connecting side path has fixed endpoints.
All spin-two cap-edge variations require a moving longitudinal endpoint.
Since the action is a sum of hinge-area and incident-simplex angle contributions, its unreduced mixed cap second derivative is therefore zero.
The integers $k_h$ are locally constant on the assumed branch, so their area terms have the same support and introduce no extra mixed derivative.
Only elimination of the single internal spin-two pair contributes to that block in Eq.~\eqref{eq:schur}.
This gives Eq.~\eqref{eq:bottleneck}.
An outer product has rank one precisely when both vectors are nonzero; tensoring with $\Id_2$ doubles that rank.
\end{proof}

\begin{remark}
The theorem explains persistence of a rank-two carrier as an architectural bottleneck.
The nonzero couplings and the curvature response remain dynamical outputs.
Adding another internal spin-two copy, breaking the symmetry or introducing nonlocal action terms can change the bound.
Rank persistence in this cone is not independent evidence for a universal continuum two-helicity theory.
\end{remark}

\section{Canonical reduction and endpoint identification}
\label{sec:canonical}
For a nonsingular scalar mixed block, write the one-polarization principal function as
\begin{equation}
S(x,y)=\tfrac12ax^2+bxy+\tfrac12cy^2,
\quad p=-\partial_xS,\quad P=\partial_yS.
\label{eq:principal}
\end{equation}
Solving for $(y,P)$ gives
\begin{equation}
\begin{pmatrix}y\\P\end{pmatrix}
=M\begin{pmatrix}x\\p\end{pmatrix},\qquad
M=\begin{pmatrix}-a/b&-1/b\\b-ca/b&-c/b\end{pmatrix}.
\label{eq:canonical}
\end{equation}
For $\Omega=\left(\begin{smallmatrix}0&1\\-1&0\end{smallmatrix}\right)$,
\begin{equation}
M^T\Omega M=\Omega,\qquad
\det(\lambda\Id-M)=\lambda^2+\frac{a+c}{b}\lambda+1.
\label{eq:spectrum}
\end{equation}
Each identity holds twice for the two scalar polarization blocks.
They follow from a nondegenerate quadratic generating function; their gravitational meaning depends on the actual Regge construction in Section~\ref{sec:action}.

At $m>1$, $B_m$ is singular on the full cap space.
The retained producer selects its two leading singular-vector axes and restricts the self blocks to those axes.
It does not invert the mixed block on the remaining detail coordinates.
To state this reduction explicitly, split initial variables into survivor $x$ and detail $d$, and final variables into $y,D$.
Writing the initial self block as $A$, the initial detail pre-constraint is
\begin{equation}
C_d=p_d+A_{ds}x+A_{dd}d=0.
\label{eq:preconstraint}
\end{equation}
The survivor momentum $\pi=p_s+A_{sd}d$ commutes with $C_d$.
At the final cap, the analogous momentum is $\Pi=P_s-C_{sd}D$.
The restricted principal relation uses $(x,\pi)$ and $(y,\Pi)$, or equivalently the section $d=D=0$.
The self-block survivor/detail coupling need not vanish; its norm is about $1.42048$ at the fifth level.
Further gluing can impose additional constraints, so this reduction does not define a universal propagator for arbitrary future extensions~\cite{Hoehn2014}.

\begin{proposition}[A tracked endpoint identification]
\label{prop:endpoints}
Let $M:V_i\to V_f$ be symplectic between two phase spaces.
Under independent symplectic coordinate changes $C_i,C_f$, its matrix becomes
\begin{equation}
M'=C_fMC_i^{-1}.
\label{eq:left_right}
\end{equation}
Its spectrum need not be invariant.
If a symplectic identification $\mathcal J:V_i\to V_f$ is specified and transformed as
$\mathcal J'=C_f\mathcal JC_i^{-1}$, then the return map
\begin{equation}
\mathcal R=\mathcal J^{-1}M
\quad\text{satisfies}\quad
\mathcal R'=C_i\mathcal RC_i^{-1}.
\label{eq:returnmap}
\end{equation}
Its spectrum is invariant.
\end{proposition}
\begin{proof}
Equation~\eqref{eq:left_right} follows by substituting the two coordinate changes into $z_f=Mz_i$.
Inserting the transformation of $\mathcal J$ gives Eq.~\eqref{eq:returnmap}.
For a counterexample to invariance without this identification, take
\begin{equation}
M=\begin{pmatrix}0&-1\\1&0\end{pmatrix},\quad
C_i=\Id,\quad C_f=\begin{pmatrix}1&0\\-3&1\end{pmatrix}.
\label{eq:counterexample}
\end{equation}
These matrices are symplectic, but $\tr M=0$ and $\tr(C_fM)=3$.
The displayed spectrum changes from elliptic to hyperbolic.
\end{proof}

For the numerical matrices below, the chosen reduced endpoint coordinates are identified componentwise, $\mathcal J=\Id$.
Their spectra are properties of this convention.
A geometric identification motivated by time reflection must include the cotangent and orientation conventions and must be transported through the survivor projections.
That complete refinement transport has not been independently verified here.
Only common-frame transformations, or endpoint transformations that also transport $\mathcal J$, license the spectral no-rescue statement.
Common scalar coordinate dilation likewise leaves $a/b$ unchanged.

For specified coarse-to-fine phase-space maps $I_i,I_f$, a more complete consistency question is
\begin{equation}
M_{\mathrm f}I_i=I_fM_{\mathrm c},
\label{eq:intertwiner}
\end{equation}
with the constraint reduction included.
Vanishing or nonvanishing of its defect is invariant when all maps transform together.
Its numerical norm additionally needs chosen metrics.
Matching eigenvalues alone is weaker than verifying Eq.~\eqref{eq:intertwiner}.

\section{Finite dynamics and partial subdivision}
\label{sec:numerics}
The coarse time-symmetric retained computation gives, in its normalized pair bases,
\begin{equation}
A\simeq-0.57611003\Id_2,\qquad
B\simeq0.84000291\Id_2,\qquad
C\simeq-0.57611004\Id_2,
\label{eq:blocks}
\end{equation}
with imaginary and nonscalar differences of order $10^{-7}$.
The two mixed singular values are $0.84000300195$ and $0.84000282365$.
Its two positive eigenvalue arguments are $0.81503482086$ and $0.81503524166$, with maximum modulus defect $2.12\times10^{-7}$.
This is a finite elliptic canonical channel to differentiation accuracy.
The complete-action replay producing these outputs took $12.70$ seconds in the recovery calculation; this paper's default replay reconstructs the transfer from the saved blocks.

For real time-symmetric scalar blocks put $c=a$ and $\chi=-a/b$.
Then
\begin{equation}
M=\begin{pmatrix}\chi&-1/b\\b(1-\chi^2)&\chi\end{pmatrix},
\qquad\lambda^2-2\chi\lambda+1=0.
\label{eq:chi}
\end{equation}
Thus $|\chi|<1$ is elliptic, $|\chi|>1$ hyperbolic; equality is the parabolic boundary.
Table~\ref{tab:flow} combines the retained earlier levels with the completed first- and fifth-level reproductions.
The rank threshold in the numerical producer is a relative mixed singular-value cutoff $10^{-5}$; the analytic bound is Theorem~\ref{thm:rank}, not a floating-point rank declaration.

\begin{table}[htbp]
\centering\small
\caption{The selected partial-refinement channel. Levels 2--4 are retained historical complete-action values, not newly recomputed for this paper. Levels 1 and 5 have retained completed-recovery full-producer outputs. All spectra use the stated reduced endpoint identification.}
\label{tab:flow}
\begin{tabular}{rrrrl}
\toprule
$m$&leading mixed singular&largest detail singular&coarse overlap&spectrum descriptor\\
\midrule
1&0.840003&---&1.000000&$\theta\simeq0.815035$\\
2&0.672133&$3.4\times10^{-8}$&0.999903&$\theta\simeq0.832011$\\
3&0.573840&$6.6\times10^{-8}$&0.970068&$\theta\simeq0.634608$\\
4&0.516755&$8.6\times10^{-8}$&0.931103&$\theta\simeq0.388070$\\
5&0.479964&$1.85\times10^{-7}$&0.891918&$\lambda\simeq1.29025,0.77504$\\
\bottomrule
\end{tabular}
\end{table}

The coarse injection used to define the overlap is the piecewise-affine cap sample vector
\begin{equation}
I_m^{\mathrm{sample}}=
\begin{pmatrix}m^{-1}\Id_2\\2m^{-1}\Id_2\\\vdots\\\Id_2\end{pmatrix}.
\label{eq:injection}
\end{equation}
It is transformed into the orthonormal quotient basis and compared by singular values of the overlap between its orthonormal columns and the mixed-block survivor axes.
This is a configuration-subspace comparison, not the full cotangent intertwiner in Eq.~\eqref{eq:intertwiner}.

At the fifth level the four retained transfer eigenvalues are
\begin{align}
1.29024266-1.16\times10^{-6}i,&\quad
1.29024809-0.84\times10^{-6}i,\notag\\
0.77504801+0.70\times10^{-6}i,&\quad
0.77504474+0.51\times10^{-6}i.
\label{eq:fifth}
\end{align}
The departure from the unit circle is approximately $0.290248$, far above the displayed imaginary residuals.
The reduced symplectic residual is $1.76\times10^{-16}$ and the smallest physical internal Hessian singular value is $0.31048966$.
This full-producer calculation completed in $423.37$ seconds during the recovery and was not rerun for manuscript delivery.
An earlier second-step calculation is part of the original report, but the portable package does not claim to regenerate or independently recertify that second step.
Tiny positive arguments returned by a generic eigenphase field at this level are numerical imaginary contamination on a reciprocal-real spectrum; they are not an oscillation phase.

The finite dependence is a diagnostic of the stated family and convention.
It cannot establish asymptotic nonconvergence from five values.
Moreover, every level retains transverse edges of length exactly one, so maximum simplex diameter in the fixed background chart does not approach zero.
The shrinking Gram margin in Table~\ref{tab:geometry} reflects increasingly thin anisotropic simplices, not a full shape-regular bulk-and-boundary refinement.
Uniform partition $m'$ contains partition $m$ only when $m$ divides $m'$.
Thus $1\to2\to4$ is nested, whereas $2\to3$, $3\to4$ and $4\to5$ are not successive nesting steps, although each refines level one.
The data do not reject other Regge refinement families or a continuum limit.

\section{A positive invariant and a closed reduced calibration}
\label{sec:repair}
Two small exact results help interpret the finite output without replacing its open questions by a larger programme.

\begin{proposition}[Definite invariant in the elliptic scalar channel]
If $b\ne0$ and $|\chi|<1$ are real, the matrix in Eq.~\eqref{eq:chi} preserves
\begin{equation}
G=\diag\bigl(b^2(1-\chi^2),1\bigr)>0,
\qquad M^TGM=G.
\label{eq:positiveG}
\end{equation}
\end{proposition}
\begin{proof}
Put $s^2=1-\chi^2>0$.
Multiplication of $M^T\diag(b^2s^2,1)M$ gives diagonal entries $b^2s^2(\chi^2+s^2)$ and $\chi^2+s^2$, with zero off-diagonal entries.
These are exactly the entries of $G$.
\end{proof}

This is a definite conserved finite quadratic norm, not a uniquely selected positive physical Hamiltonian.
Let $\theta=\arccos\chi\in(0,\pi)$ and $s=\sin\theta$.
The principal logarithm is $L=(\theta/s)(M-\chi\Id)$, and in the convention $M=\exp(\Omega K)$ its quadratic generator is
\begin{equation}
K=-\Omega L=
\diag\left(-\theta b s,-\frac{\theta}{bs}\right).
\label{eq:log}
\end{equation}
For the positive $b$ in the witness, this branch is negative definite.
Replacing $\theta$ by $\theta-2\pi$ produces a positive definite generator of the same one-step map.
The map alone does not choose this branch, an oriented physical time duration, a quantum positive-frequency vacuum or a ghost-free measure.

\begin{theorem}[Matching a selected scalar kernel]
\label{thm:repair}
Suppose $S_m$ has real time-symmetric scalar blocks $a_m,b_m$, with $b_m\ne0$.
Choose target coefficients $a_*=-b_*\cos\theta_*$ with $b_*/b_m>0$.
There is a unique common self-block correction at fixed $b_m$ that matches the target transfer trace:
\begin{equation}
\Delta S_m=\frac{\delta a_m}{2}(|x|^2+|y|^2),
\qquad\delta a_m=-b_m\cos\theta_*-a_m.
\label{eq:counterterm}
\end{equation}
After the common substitution $x=r_mx_*$, $y=r_my_*$ with $r_m^2=b_*/b_m$, the entire corrected scalar two-polarization kernel equals the target kernel.
\end{theorem}
\begin{proof}
The corrected trace is $-2(a_m+\delta a_m)/b_m$.
Setting it to $2\cos\theta_*$ gives Eq.~\eqref{eq:counterterm} uniquely.
The substitution multiplies both coefficients by $r_m^2$.
It sends $b_m$ to $b_*$ and the corrected $a_m$ to $-b_*\cos\theta_*=a_*$.
Both self terms and the mixed term therefore agree, for each component.
\end{proof}

For $\theta_*=0.8150346718577313$ and $b_*=0.840003$, the rounded levels 2--4 in Table~\ref{tab:flow} require
\begin{equation}
\delta a_2\simeq-0.00836990,\quad
\delta a_3\simeq0.06855150,\quad
\delta a_4\simeq0.12391678.
\label{eq:repairs}
\end{equation}
The real scalar projection of the retained fifth-level blocks, defined by
$a_5=\Re\tr(A+C)/4$ and $b_5=\Re\tr B/2$, gives
\begin{align}
a_5&=-0.49563313524313946,&b_5&=0.4799642815824481,\notag\\
\chi_5&=1.032645874415969,&\delta a_5&=0.16645291498144332,\notag\\
r_5&=1.322927240460913.
\label{eq:repair5}
\end{align}
The theorem exactly repairs this real scalar projection, while the original complex numerical block retains its recorded errors.

This construction changes the selected action; it does not contradict the common-coordinate phase invariance or undo the hyperbolic matrix by a conjugacy.
Here ``kernel'' denotes the quadratic principal function, without an uncomputed quantum prefactor or measure.
It also shows the limits of a trace-matching condition.
A projector onto a global survivor can be nonlocal on the fine boundary, and matching this channel supplies no proof of locality, full Regge-constraint preservation, compatible gluing, control of other sectors or a perfect four-dimensional action.
Those additional requirements cannot be inferred from Theorem~\ref{thm:repair}.

\section{Evidence, reproducibility and limits}
\label{sec:evidence}
The article separates exact algebra, computed finite data and physical interpretation.
The commutant, rank bound under its hypotheses, canonical identities, endpoint-identification proposition, positive invariant and scalar repair have explicit finite proofs.
The simplex signatures and boundary classifications are numerical checks on explicit input coordinates.
The first- and fifth-level complete-action outputs were replayed in the preceding recovery; their input/output bytes and producer dependencies are frozen in the package with SHA-256 provenance.
Levels 2--4 and the original second-step sensitivity checks remain inherited evidence, not new executions for this manuscript.
No independently implemented complete Regge action has been used to reproduce the numerical Hessian.

The default command \texttt{python reproduce.py} needs only packaged inputs and ordinary Python dependencies.
It verifies the frozen-file hashes, solves the exact algebra, rebuilds the geometry, checks the architectural support condition, reconstructs the first- and fifth-level transfer matrices from their saved blocks, and verifies their spectra and scalar repairs.
It does not recompute the Regge Hessian.
The separate optional \texttt{run\_full.py} executes the frozen original producer on level 1 or 5; there is no hidden default expensive run or dependency on an old local project path.

The finite computations support a causal dihedral angular-weight-two transfer channel, with an elliptic coarse witness and a convention-specific fifth-level hyperbolic transfer.
The geometry gives a time-oriented region, not a retarded Green function.
Discrete angular sampling, finite-difference residuals, constrained reduction and endpoint identification all restrict the interpretation.
A physically selected positive-energy branch, massless pole, continuum lightlike dispersion, nonlinear stability, quantum measure and shape-regular continuum refinement have not been established.
These gaps do not erase the explicit finite result; they delimit its question.

The integrated causal join-star benchmark is the contribution offered for further assessment.
The classical ingredients credited above, elementary outer-product and canonical algebra, and the reduced counterterm fit are not claimed as new general theories.
The primary-source comparison was targeted; it supplies no comprehensive priority certificate.

\section{Conclusion}
The explicit Lorentzian cone supports a nonzero finite dihedral spin-two transfer pair.
Its mixed rank is controlled by a single internal angular channel, rather than by an independent continuum helicity count.
Partial longitudinal subdivision changes the selected finite dynamics and reaches a hyperbolic fifth-level transfer, but leaves large transverse edges and does not test a full continuum limit.
Endpoint spectra become invariant only after fixing and tracking an identification of the endpoint phase spaces.
The resulting finite scalar channel admits an explicit calibrated kernel repair whose full-lattice realization remains a separate question.

\section*{Contribution and review disclosure}
Human authors and affiliations are intentionally blank, and human expert review has not occurred.
The operator reports that earlier exploratory work was led by Fable 5, while results from July 2026 onward were predominantly produced by GPT models; DeepSeek and multiple GPT models also participated in review and preparation.
This reported chronology does not allocate individual authorship, intellectual credit or ownership.
The operator's attribution statement is: ``I genuinely do not know how authorship, intellectual credit, or ownership should be attributed for these results.''

\bibliographystyle{unsrt}
\bibliography{references}
\end{document}
